\documentclass[11pt]{article}
\usepackage{amsmath,amssymb,amsthm}
\usepackage[margin=1in]{geometry}
\usepackage{booktabs}
\usepackage{hyperref}
\usepackage{microtype}

\newtheorem{theorem}{Theorem}
\newtheorem{proposition}[theorem]{Proposition}
\newtheorem{corollary}[theorem]{Corollary}
\newtheorem{lemma}[theorem]{Lemma}
\newtheorem{definition}[theorem]{Definition}
\newtheorem{remark}[theorem]{Remark}

\newcommand{\Z}{\mathbb{Z}}
\newcommand{\R}{\mathbb{R}}
\newcommand{\N}{\mathbb{N}}
\newcommand{\Caso}{C}
\newcommand{\CasoM}{W}
\DeclareMathOperator{\sgn}{sgn}

\title{A machine-checked, uniform-in-order Abel formula for the Casoratian,\\
with a formalized fundamental-system converse and an irrationality certificate,\\
in Lean~4/Mathlib}
\author{Papanokechi\\
\small ORCID: \href{https://orcid.org/0009-0000-6192-8273}{0009-0000-6192-8273}}
\date{June 2026 \quad(\texttt{pcf-casoratian-catalogue} v1.1)}

\begin{document}
\maketitle

\begin{abstract}
For $k$ sequences $s^{(0)},\dots,s^{(k-1)}\colon\N\to R$ over a commutative ring $R$,
all solving the order-$k$ homogeneous linear recurrence
$s(n+k)=\sum_{j=0}^{k-1}c_j(n)\,s(n+j)$, the \emph{Casoratian} (discrete Wronskian)
$\Caso(n)=\det\big[s^{(t)}(n+i)\big]_{0\le i,t<k}$ obeys the first-order law
$\Caso(n+1)=(-1)^{k-1}c_0(n)\,\Caso(n)$, hence the closed product form
$\Caso(n)=\big(\prod_{l<n}(-1)^{k-1}c_0(l)\big)\Caso(0)$. Only the \emph{lowest}
coefficient $c_0$ survives; $c_1,\dots,c_{k-1}$ cancel. This is \emph{Abel's formula}
for the Casoratian --- classical, going back to Casorati, the discrete analogue of
Liouville's formula for the Wronskian. Our contribution is not the identity but its treatment: we give a single
proof valid \emph{uniformly in the order} $k$ and formalize it in Lean~4/Mathlib to a
clean axiom cone (a subset of $\{\texttt{propext},\texttt{Classical.choice},
\texttt{Quot.sound}\}$, no \texttt{sorry}). The formal proof rewrites the last row of
$\Caso(n+1)$, via the recurrence, into a linear combination of the
\emph{cyclically rotated} rows of $\Caso(n)$; \texttt{Matrix.det\_updateRow\_sum}
collapses the combination to the single coefficient $c_0(n)$ and
\texttt{sign\_finRotate} supplies the cyclic sign $(-1)^{k-1}$. We also formalize the
constant-coefficient power form $\Caso(n)=\big((-1)^{k-1}c_0\big)^n\Caso(0)$, an
integral-domain non-vanishing (linear-independence) certificate, and the
order-$2$ and order-$3$ instances proved independently by hand as faithfulness
witnesses (order $2$ reproduces the corpus sign $-c_0$, with Euler's $(-1)^n$
convergent Wronskian as the case $c_0\equiv1$; order $3$ gives $+c_0$). The result
opens the order-$\ge3$ direction in a formalized polynomial-continued-fraction corpus
previously confined to the order-$2$ (three-term) case --- the algebraic backbone of
simultaneous / Hermite--Pad\'e approximation and of Ap\'ery-type irrationality
certificates. Building on the law we further formalize its \emph{fundamental-system
converse} --- forward determinacy of a solution by its first $k$ values, and
completeness of the basis once the initial Casoratian is a unit --- so that the
solutions form a machine-checked \emph{fundamental set}; and we put that set to work
through a reusable, machine-checked \emph{irrationality pipeline} that turns any integer
fundamental set together with a shrinking \emph{real} Hermite--Pad\'e remainder into a
proof of irrationality, the Casoratian discharging the non-vanishing
(linear-independence) input \emph{structurally}. We instantiate it end-to-end on the
order-$3$ recurrence $u(n{+}3)=2u(n{+}2)-u(n)$ (characteristic roots $\varphi,\psi,1$),
whose fundamental set $\{1,F_n,F_{n+1}\}$ and remainder $F_n\varphi-F_{n+1}=-\psi^n\to0$
re-derive the irrationality of the golden ratio $\varphi$ through the Casoratian --- an
independent route from Mathlib's $\sqrt5$-based proof, and a demonstration that the
determinantal machinery genuinely closes an irrationality theorem. We claim no novelty
for $\varphi$'s irrationality itself, which is classical; the contributions are the
uniform-in-$k$ machine-checked law, the formalized converse, and the reusable
certificate engine. An independent computer-algebra gate confirms the identity
symbolically at $k=3$, by exact-rational checks for $k=2,\dots,7$, and on named
instances.
\end{abstract}

% ---------------------------------------------------------------------------
% SHORT ABSTRACT (<=210 words) -- drop-in variant for venues that cap the
% abstract length.  Swap the abstract environment above for this block.
%
% \begin{abstract}
% For $k$ sequences over a commutative ring solving a common order-$k$ linear
% recurrence, the Casoratian (discrete Wronskian) $C(n)=\det[s^{(t)}(n+i)]$ obeys
% Abel's first-order law $C(n+1)=(-1)^{k-1}c_0(n)\,C(n)$, so only the lowest
% coefficient survives.  The identity is classical; our contribution is its
% treatment.  We give a single proof valid \emph{uniformly in the order} $k$ and
% formalize it in Lean~4/Mathlib to a clean axiom cone, together with the closed
% product form, the constant-coefficient form, and an integral-domain
% non-vanishing (linear-independence) certificate.  Building on it we formalize the
% \emph{fundamental-system converse} -- forward determinacy and completeness once
% the initial Casoratian is a unit -- and a reusable, machine-checked
% \emph{irrationality pipeline} that turns any integer fundamental set plus a
% shrinking real Hermite--Pad\'e remainder into a proof of irrationality, the
% Casoratian discharging the non-vanishing input structurally.  We instantiate it
% on an order-$3$ recurrence to re-derive the irrationality of the golden ratio
% through the Casoratian, claiming no novelty for $\varphi$ itself.  An independent
% computer-algebra gate confirms the law symbolically.  Twenty-six declarations are
% machine-checked.
% \end{abstract}
% ---------------------------------------------------------------------------

\noindent\textbf{2020 MSC:} 39A06 (primary), 68V20, 11J72, 15A15.\quad
\textbf{Keywords:} Casoratian, discrete Wronskian, Abel's formula,
linear difference equation, higher-order recurrence, fundamental system,
Hermite--Pad\'e approximation, Ap\'ery-type irrationality, irrationality criterion,
golden ratio, Lean~4, Mathlib, formal verification, polynomial continued
fractions.

\section{Introduction}\label{sec:intro}

The convergents of a polynomial continued fraction $K_{n\ge1}\,a_n/b_n$ are governed
by a three-term recurrence $q_n=b_nq_{n-1}+a_nq_{n-2}$, and the basic structural
invariant is the order-$2$ Casoratian (discrete Wronskian) of two solutions: it
flips sign and scales by the lower coefficient at each step, so the convergent
Wronskian $p_{n+1}q_n-p_nq_{n+1}=(-1)^n$ is constant up to sign. This is the
certificate behind Euler's irrationality criterion, and across the deposited corpus
the Casoratian appears only in this order-$2$, three-term form
(\cite{ladder,caso2lean}; formalized as \texttt{casoratian\_eqG} in the companion
\texttt{pcf-delta} project~\cite{pcfdelta}).

\paragraph{The order-$\ge3$ direction.}
The natural next object is the Casoratian of $k\ge3$ solutions of an order-$k$
recurrence
\begin{equation}\label{eq:rec}
  s(n+k)=\sum_{j=0}^{k-1}c_j(n)\,s(n+j),
  \qquad c_0,\dots,c_{k-1}\colon\N\to R,
\end{equation}
the $k\times k$ determinant
\begin{equation}\label{eq:caso}
  \Caso(n)=\det\big[s^{(t)}(n+i)\big]_{0\le i,t<k}.
\end{equation}
This is the algebraic backbone of \emph{simultaneous} (Hermite--Pad\'e)
approximation, the engine of Ap\'ery-type irrationality proofs
(\cite{vdpoorten}): there one builds several sequences satisfying one common
higher-order recurrence and certifies their linear independence through the
non-vanishing of exactly this determinant. The first-order law that $\Caso$ obeys
is classical --- the discrete analogue of Abel's identity for the Wronskian of a
linear ODE, going back to Casorati; see, e.g., \cite[Ch.~2]{elaydi}
and \cite{kelleypeterson}. It states that $\Caso$ is reproduced at each step by a
single scalar built only from the \emph{lowest} coefficient $c_0$.

\paragraph{Contribution.}
The identity itself is not new. What we provide is:
\begin{enumerate}
\item[(i)] a single proof valid \emph{uniformly in the order} $k$
(Theorem~\ref{thm:ajl}), and its formalization in Lean~4 with Mathlib to a clean
axiom cone (Section~\ref{sec:lean}), so the result is \textsc{proven} in the strict
sense of this program: \texttt{\#print axioms} reports a subset of
$\{\texttt{propext},\texttt{Classical.choice},\texttt{Quot.sound}\}$ with no
\texttt{sorryAx} and no \texttt{sorry}/\texttt{admit}/\texttt{native\_decide} in
source;
\item[(ii)] the closed product form, the constant-coefficient power form, and an
integral-domain non-vanishing (linear-independence) certificate, all formalized
(Section~\ref{sec:cons});
\item[(iii)] the order-$2$ and order-$3$ instances proved \emph{independently} by a
direct cofactor computation, kept as faithfulness witnesses against the general
proof and machine-checked to coincide with the general determinant
(Proposition~\ref{prop:lowk}); and
\item[(iv)] an independent computer-algebra gate (Section~\ref{sec:gate}) that
re-derives the law symbolically and on named instances, in the program's
``gate-before-Lean'' discipline;
\item[(v)] the \emph{fundamental-system converse} (Section~\ref{sec:converse}),
formalized: forward determinacy of a solution by its first $k$ values and completeness
of the basis once the initial Casoratian is a unit, so that the solutions form a
machine-checked fundamental set --- the object Hermite--Pad\'e / Ap\'ery applications
consume; and
\item[(vi)] a reusable, machine-checked \emph{irrationality pipeline}
(Section~\ref{sec:cert}) that converts any integer fundamental set plus a shrinking
real Hermite--Pad\'e remainder into a proof of irrationality --- with the non-vanishing
input discharged structurally by the Casoratian --- instantiated end-to-end on a
concrete order-$3$ recurrence to re-derive the irrationality of the golden ratio
$\varphi$ through the discrete Wronskian.
\end{enumerate}
Concretely, this opens the order-$\ge3$ space inside a formalized PCF corpus that had
treated only order $2$, supplies a reusable, machine-checked independence
certificate for future Hermite--Pad\'e constructions, and --- through the converse and
the irrationality pipeline --- carries that certificate all the way to a worked
irrationality theorem.

\section{The Casoratian of a linear recurrence}\label{sec:setup}

Throughout, $R$ is a commutative ring and $k\ge1$ a fixed order. We write a solution
of~\eqref{eq:rec} simply as a sequence $s\colon\N\to R$ satisfying that identity for
all $n$; no invertibility of the coefficients is assumed.

\begin{definition}[Casoratian matrix and Casoratian]\label{def:caso}
For sequences $s^{(0)},\dots,s^{(k-1)}\colon\N\to R$, the \emph{Casoratian matrix}
$\CasoM(n)$ at $n$ is the $k\times k$ matrix whose $(i,t)$ entry is the $i$-shift of the
$t$-th sequence, $s^{(t)}(n+i)$. The \emph{Casoratian} is its determinant,
$\Caso(n)=\det\CasoM(n)$, the scalar of~\eqref{eq:caso}. Throughout we write $\CasoM$
for the matrix and reserve $\Caso=\det\CasoM$ for the scalar determinant; in the
formalization the matrix is \texttt{casoMat} and the Casoratian is
\texttt{(casoMat\,s\,n).det}.
\end{definition}

Because $\det$ is invariant under transposition, it is immaterial whether one indexes
rows by shifts and columns by solutions or vice versa; we fix rows~$=$~shifts. The
two ``ends'' of the next matrix $\CasoM(n+1)$ are what drive the law:
\begin{itemize}
\item \textbf{Top $k-1$ rows.} For $0\le i\le k-2$, row $i$ of $\CasoM(n+1)$ is the
shift vector $\big(s^{(t)}(n+1+i)\big)_t$, which is precisely row $i+1$ of
$\CasoM(n)$. So the top $k-1$ rows of $\CasoM(n+1)$ are the bottom $k-1$ rows of
$\CasoM(n)$, in order.
\item \textbf{Bottom row.} Row $k-1$ of $\CasoM(n+1)$ is the shift vector at
$n+k$, $\big(s^{(t)}(n+k)\big)_t$, which by the recurrence~\eqref{eq:rec} equals
$\sum_{j=0}^{k-1}c_j(n)\,\big(s^{(t)}(n+j)\big)_t$, a linear combination of \emph{all}
$k$ rows of $\CasoM(n)$.
\end{itemize}

\section{The general-order Abel formula for the Casoratian}\label{sec:law}

\begin{theorem}[Abel's formula for the Casoratian, uniform in $k$]\label{thm:ajl}
Let $s^{(0)},\dots,s^{(k-1)}\colon\N\to R$ all satisfy the order-$k$
recurrence~\eqref{eq:rec}. Then the Casoratian~\eqref{eq:caso} satisfies the
first-order law
\begin{equation}\label{eq:step}
  \Caso(n+1)=(-1)^{k-1}\,c_0(n)\,\Caso(n)\qquad(n\in\N),
\end{equation}
and consequently the closed product form
\begin{equation}\label{eq:prod}
  \Caso(n)=\Bigg(\prod_{l=0}^{n-1}(-1)^{k-1}c_0(l)\Bigg)\Caso(0)
          =(-1)^{(k-1)n}\Bigg(\prod_{l=0}^{n-1}c_0(l)\Bigg)\Caso(0).
\end{equation}
Only the lowest coefficient $c_0$ enters; $c_1,\dots,c_{k-1}$ cancel.
\end{theorem}

\begin{proof}
Let $\rho$ be the cyclic permutation of $\{0,1,\dots,k-1\}$ that moves every index up
by one and wraps the top to the bottom, $\rho(i)=i+1$ for $i<k-1$ and
$\rho(k-1)=0$; it is a single $k$-cycle, so $\sgn\rho=(-1)^{k-1}$. Let $M$ be
$\CasoM(n)$ with its rows permuted by $\rho$, i.e.\ row $i$ of $M$ is row $\rho(i)$ of
$\CasoM(n)$. By the observation in Section~\ref{sec:setup}, the top $k-1$ rows of $M$
(namely rows $1,\dots,k-1$ of $\CasoM(n)$) are exactly the top $k-1$ rows of
$\CasoM(n+1)$, while the bottom row of $M$ is row $0$ of $\CasoM(n)$.

Replacing the bottom row of $M$ by the linear combination
$\sum_{j=0}^{k-1}c_j(n)\,\big(\text{row }j\text{ of }\CasoM(n)\big)$ produces exactly
$\CasoM(n+1)$, because that combination is the bottom row of $\CasoM(n+1)$ and the top
rows already agree. Now take determinants. Multilinearity in the replaced row --- the
content of \texttt{Matrix.det\_updateRow\_sum} --- gives
\[
  \det\CasoM(n+1)
   =\sum_{j=0}^{k-1}c_j(n)\,\det\!\big(M\text{ with bottom row }=\text{row }j\text{ of }\CasoM(n)\big).
\]
Every summand with $j\ge1$ places into the bottom position a row of $\CasoM(n)$ that
already occurs among the top $k-1$ rows of $M$, so that determinant vanishes. Only
$j=0$ survives, and there the bottom row is the original row $0$ of $\CasoM(n)$, so the
summand is $c_0(n)\,\det M$. Finally $\det M=\sgn\rho\cdot\det\CasoM(n)
=(-1)^{k-1}\det\CasoM(n)$. Recalling $\Caso=\det\CasoM$, combining gives~\eqref{eq:step}.
A trivial induction on $n$, using $\Caso(0)$ as the base case and~\eqref{eq:step} at the
inductive step, yields the product form~\eqref{eq:prod}.
\end{proof}

\begin{remark}[What the formalization does, verbatim]\label{rem:lean-proof}
The Lean development (Section~\ref{sec:lean}) realizes this proof literally. It builds
$M$ as \texttt{(casoMat\,s\,n).submatrix (finRotate (k+1)) id} (Mathlib's
\texttt{finRotate} is the cycle $\rho$), expresses $\CasoM(n+1)$ as that submatrix with
its last row updated to the recurrence combination, then rewrites with
\texttt{Matrix.det\_updateRow\_sum} (the multilinear collapse to the surviving
coefficient) and \texttt{Matrix.det\_permute} together with \texttt{sign\_finRotate}
(the cyclic sign). The ``vanishing for $j\ge1$'' step is automatic: it is exactly the
statement that \texttt{det\_updateRow\_sum} isolates the coefficient indexed by the
updated row. No analytic input is used; the argument is finitary algebra over an
arbitrary \texttt{CommRing}.
\end{remark}

\section{Specializations and consequences}\label{sec:cons}

\begin{corollary}[Constant lowest coefficient]\label{cor:const}
If $c_0(l)=a$ for all $l$, then $\Caso(n)=\big((-1)^{k-1}a\big)^n\Caso(0)$. In
particular, if $a$ is a unit (or $1$), the Casoratian never vanishes once
$\Caso(0)\ne0$.
\end{corollary}

\begin{corollary}[Non-vanishing / linear-independence certificate]\label{cor:indep}
Suppose $R$ is an integral domain. If $\Caso(0)\ne0$ and $c_0(l)\ne0$ for every
$l<n$, then $\Caso(n)\ne0$. Hence if $\Caso(0)\ne0$ and $c_0$ is nowhere zero, the
$k$ solutions stay linearly independent at every window: $\Caso(n)\ne0$ for all $n$.
\end{corollary}

\begin{proof}
Immediate from~\eqref{eq:prod}: a product of nonzero elements of a domain is nonzero,
and $(-1)^{k-1}$ is a unit.
\end{proof}

Corollary~\ref{cor:indep} is the order-$k$ analogue of the order-$2$ statement
$p_{n+1}q_n-p_nq_{n+1}=(-1)^n\ne0$ that underlies Euler's criterion: it is precisely
the determinantal independence certificate one needs for Hermite--Pad\'e systems.

\begin{proposition}[Low-order witnesses]\label{prop:lowk}
The cases $k=2,3$ hold and are proved independently by direct cofactor expansion,
without the permutation argument:
\[
  k=2:\quad \Caso(n+1)=-\,c_0(n)\,\Caso(n);
  \qquad
  k=3:\quad \Caso(n+1)=+\,c_0(n)\,\Caso(n).
\]
The order-$2$ sign $-c_0$ is the corpus sign; its $c_0\equiv1$ instance is the
constant convergent Wronskian $p_{n+1}q_n-p_nq_{n+1}=(-1)^n$.
\end{proposition}

These hand proofs (\texttt{caso2\_step}, \texttt{caso3\_step}; the $3\times3$
cofactor expansion is closed by \texttt{ring}) are retained as \emph{faithfulness
witnesses}: they agree with Theorem~\ref{thm:ajl} at $k=2,3$, guarding against a
vacuous or mis-signed general statement. Moreover, the bespoke order-$3$ expansion is
reconciled with the general object by the machine-checked bridge
\texttt{caso3\_matches\_general}, which proves \texttt{caso3 y z w n} equal to
$\det\CasoM(n)$ for the general Casoratian matrix \texttt{casoMat} (Section~\ref{sec:law})
specialized to $k=2$ with $\{y,z,w\}$; so the witnesses certify the \emph{formalized}
general theorem (Theorem~\ref{thm:lean}), not merely a parallel polynomial identity.

\section{The fundamental-system converse}\label{sec:converse}

The Abel law and its non-vanishing corollary say that the chosen solutions stay
\emph{independent}; the converse direction --- that they are also \emph{complete},
i.e.\ span every solution --- is what upgrades $\{s^{(t)}\}$ from an independent family
to a genuine \emph{fundamental set}, the object Hermite--Pad\'e and Ap\'ery
constructions actually consume. We formalize it in two parts.

\begin{theorem}[Fundamental-system converse; \textsc{proven}]\label{thm:converse}
Let $s^{(0)},\dots,s^{(k-1)}\colon\N\to R$ solve the order-$k$ recurrence~\eqref{eq:rec}.
\begin{enumerate}
\item[\textup{(a)}] \emph{(Forward determinacy.)} Any solution of~\eqref{eq:rec} is
determined by its first $k$ values: if $y,z\colon\N\to R$ both satisfy~\eqref{eq:rec}
and $y(i)=z(i)$ for $0\le i<k$, then $y(n)=z(n)$ for all $n$. No invertibility
hypothesis is needed.
\item[\textup{(b)}] \emph{(Completeness.)} If the initial Casoratian $\Caso(0)$ is a
unit of $R$, then every solution $y$ of~\eqref{eq:rec} is an $R$-linear combination of
the $s^{(t)}$: there is $a\colon\{0,\dots,k-1\}\to R$ with
$y(n)=\sum_{t} a_t\, s^{(t)}(n)$ for all $n$, given explicitly by the coordinates
$a=\CasoM(0)^{-1}\,\big(y(0),\dots,y(k-1)\big)^{\!\top}$.
\end{enumerate}
Together with the non-vanishing certificate (Corollary~\ref{cor:indep}), $\{s^{(t)}\}$
is then a machine-checked fundamental set.
\end{theorem}

\begin{proof}[Proof (as formalized)]
Part (a) is a strong induction on $n$: for $n<k$ the hypothesis gives $y(n)=z(n)$
directly, and for $n\ge k$ the recurrence writes both $y(n)$ and $z(n)$ as the
\emph{same} $R$-combination of strictly earlier values, equal by the inductive
hypothesis. For part (b), set $a=\CasoM(0)^{-1}\,(y(i))_{i<k}$ --- legitimate because
$\Caso(0)$ is a unit, so the matrix $\CasoM(0)$ is invertible --- so that
$\sum_t a_t\,s^{(t)}(i)=y(i)$ on the first window $0\le i<k$. The combination
$\sum_t a_t\,s^{(t)}$ is itself a solution of~\eqref{eq:rec}, so by part (a) it agrees
with $y$ at \emph{every} $n$. In particular completeness needs only that $\Caso(0)$ is a
unit: the forward determinacy of (a) carries the initial-window agreement to all $n$
with no reversibility --- no hypothesis on $c_0$ --- required, a small simplification
over the two-hypothesis form in which we previously recorded this converse as merely
\textsc{structural}.
\end{proof}

\begin{remark}[Lean]\label{rem:converse-lean}
Part (a) is \texttt{PcfGeneralCaso.sol\_unique\_of\_init} (strong induction,
\texttt{Nat.strong\_induction\_on}); part (b) is
\texttt{PcfGeneralCaso.solution\_isLinearCombo}, which takes coordinates from Mathlib's
nonsingular inverse (\texttt{Matrix.mul\_nonsing\_inv} at the initial window) and then
invokes part (a) to propagate. Both are in \texttt{GeneralCaso.lean} (order $k+1$ via
\texttt{Fin (k+1)}, Remark~\ref{rem:index}), with cone
$\{\texttt{propext},\texttt{Classical.choice},\texttt{Quot.sound}\}$.
\end{remark}

\section{An irrationality certificate from the fundamental set}\label{sec:cert}

We now put the fundamental set to number-theoretic work. The construction is the
discrete heart of the Hermite--Pad\'e / Ap\'ery mechanism, reduced to its
determinantal skeleton and machine-checked end to end.

\paragraph{The window is a Casoratian matrix--vector product.}
For coefficients $a\colon\{0,\dots,k-1\}\to R$, the linear form
$L(n)=\sum_t a_t\,s^{(t)}(n)$ evaluated on the consecutive window $n,\dots,n+k-1$ is
exactly $\CasoM(n)\,a$,
\begin{equation}\label{eq:window}
  \big(L(n+i)\big)_{0\le i<k}=\CasoM(n)\,a
\end{equation}
(\texttt{caso\_mulVec\_apply}). Hence, over an integral domain, if $\Caso(n)\ne0$ the
matrix $\CasoM(n)$ is injective and a \emph{nontrivial} form ($a\ne0$) cannot vanish on
a full window of $k$ consecutive indices (\texttt{linForm\_window\_ne\_zero});
consequently such a form is \emph{frequently nonzero} as $n\to\infty$
(\texttt{linForm\_frequently\_ne\_zero}). This is precisely the non-degeneracy the
irrationality machinery needs --- and the non-vanishing Casoratian supplies it for free.

\paragraph{The criterion.}
The standard linear-form irrationality criterion is
\texttt{irrational\_of\_linearForm\_frequently}: for $\alpha\in\R$ and integer sequences
$p,q$, if the form $q(n)\alpha-p(n)$ is frequently nonzero and tends to $0$, then
$\alpha$ is irrational. Were $\alpha=a/b$ rational ($b\ne0$), then
$b\,(q(n)\alpha-p(n))=q(n)a-b\,p(n)\in\Z$ is a nonzero integer whenever the form is
nonzero, so $|q(n)\alpha-p(n)|\ge1/|b|>0$ frequently, contradicting convergence to $0$.

\paragraph{The capstone.}
\texttt{irrational\_of\_fundamentalSet\_form} combines the two. Given an integer
fundamental set $s$ (so $\Caso(n)\ne0$ for all $n$) and a \emph{real} coefficient vector
$v\ne0$ for which the remainder
\[
  q(n)\alpha-p(n)=\sum_t v_t\,s^{(t)}(n)
\]
tends to $0$, the conclusion is that $\alpha$ is irrational. The frequently-nonzero
hypothesis of the criterion is discharged \emph{structurally}: the nonzero integer
Casoratian, cast into $\R$ (\texttt{Int.cast\_det} and injectivity of $\Z\hookrightarrow\R$),
stays nonzero, so by~\eqref{eq:window} the real remainder is frequently nonzero. The
only genuinely analytic input is the limit, kept as a single labelled hypothesis in the
program's established style (cf.\ the conditional limits of the corpus's Topic~3 and
Topic~5).

\begin{remark}[Why the coefficients must be real]\label{rem:realv}
The vector $v$ is real, not integer, and this is essential rather than cosmetic. With an
integer form, the remainder $q(n)\alpha-p(n)$ would be forced to be a cast integer,
which cannot shrink to $0$ while remaining frequently nonzero --- the criterion's two
hypotheses would be jointly unsatisfiable and the ``certificate'' \emph{vacuous}. The
honest Hermite--Pad\'e remainder is a genuine non-integer real combination of integer
solutions: the integer Casoratian certifies their independence, while the real
coefficients let the remainder converge.
\end{remark}

\paragraph{A concrete instance: the golden ratio.}
Take the order-$3$ recurrence
\begin{equation}\label{eq:fibrec}
  u(n+3)=2\,u(n+2)-u(n),\qquad x^3-2x^2+1=(x^2-x-1)(x-1),
\end{equation}
with characteristic roots $\varphi=\tfrac{1+\sqrt5}{2}$, $\psi=\tfrac{1-\sqrt5}{2}$,
and $1$. The three integer sequences $\{1,\,F_n,\,F_{n+1}\}$ --- the constant solution
and two shifts of Fibonacci --- solve~\eqref{eq:fibrec} (\texttt{fibSol\_rec}), and their
initial Casoratian is the Cassini determinant
\[
  \Caso(0)=\det\!\begin{pmatrix}1&0&1\\ 1&1&1\\ 1&1&2\end{pmatrix}=1
\]
(\texttt{fibCaso0}); since here $c_0\equiv-1$ is a unit, Corollary~\ref{cor:indep} gives
$\Caso(n)\ne0$ for all $n$ (\texttt{fibCaso\_ne\_zero}), so $\{1,F_n,F_{n+1}\}$ is a
fundamental set. Because $|\psi|<1$, the Hermite--Pad\'e remainder shrinks,
\[
  F_n\,\varphi-F_{n+1}=-\psi^{\,n}\longrightarrow 0
\]
(Binet's identity, Mathlib's \texttt{Real.fib\_succ\_sub\_goldenRatio\_mul\_fib};
\texttt{fib\_remainder\_tendsto}). Feeding this through the capstone with the real
coefficient vector $v=(0,\varphi,-1)$ yields
\[
  \texttt{goldenRatio\_irrational\_via\_caso}:\qquad \varphi\ \text{is irrational,}
\]
derived entirely through the discrete-Wronskian machinery.

Note that $v=(0,\varphi,-1)$ assigns coefficient $0$ to the constant solution, so the
form $F_n\varphi-F_{n+1}$ is two-term and this is deliberately the \emph{minimal}
instance --- the $2\times2$ Casoratian of $\{F_n,F_{n+1}\}$ would already certify the
non-vanishing here. What the order-$3$ pipeline genuinely exercises is its
\emph{machinery}: the $3\times3$ Casoratian matrix, its integer determinant cast into
$\R$, and the window/criterion plumbing, all run at $k=2$ (order $3$). We keep the
instance minimal on purpose, reserving a genuinely order-$3$-essential target --- one in
which all three solutions carry weight --- for the open problem of
Section~\ref{sec:status}.

\begin{remark}[Honest status of this instance]\label{rem:honest-phi}
We claim \emph{no} novelty for the irrationality of $\varphi$: it is classical and
already in Mathlib as \texttt{Real.goldenRatio\_irrational} (a $\sqrt5$-based argument).
Our \texttt{goldenRatio\_irrational\_via\_caso} is an \emph{independent route}: the
\emph{irrationality argument} does not invoke the irrationality of $\sqrt5$ --- it
performs no $\sqrt5$-descent, deriving irrationality instead from the order-$3$
Casoratian fundamental set and a shrinking remainder through the linear-form criterion.
(The roots $\varphi,\psi$ are of course \emph{defined} via $\sqrt5$, and $|\psi|<1$ is a
$\sqrt5$ fact; what differs is the \emph{mechanism} of the proof, not the avoidance of
$\sqrt5$ as a quantity.) It serves as a
\emph{faithfulness/demonstration witness} that the pipeline genuinely closes an
irrationality theorem end to end, exactly as the by-hand low-order witnesses
(Proposition~\ref{prop:lowk}) guard the general law. A geometric companion
$\{1,2^n,3^n\}$ of $u(n+3)=6u(n+2)-11u(n+1)+6u(n)$, with explicit Casoratian
$\Caso(n)=2\cdot6^n$ (\texttt{gcaso\_ne\_zero}) and window certificate
\texttt{geom\_window\_ne\_zero}, exhibits the non-vanishing half on a second recurrence
but carries \emph{no} shrinking remainder (all its roots have modulus $\ge1$), so it is
deliberately \emph{not} promoted to an irrationality certificate. What remains genuinely
open (Section~\ref{sec:status}) is to drive the same pipeline to a target \emph{not}
already known to be irrational.
\end{remark}

\section{The machine-checked formalization}\label{sec:lean}

The development is in Lean~4 with Mathlib (toolchain pinned
\texttt{leanprover/lean4:v4.30.0}, Mathlib \texttt{rev=v4.30.0}). It is built and
axiom-checked inside the companion \texttt{pcf-delta} \texttt{PcfContinuant}
project~\cite{pcfdelta}, whose \texttt{Check.lean} gate runs \texttt{\#print axioms}
on every declaration; the three source files \texttt{GeneralCaso.lean} (the
uniform-in-$k$ core and the fundamental-system converse, namespace
\texttt{PcfGeneralCaso}), \texttt{HigherCaso.lean} (the by-hand low-order witnesses,
namespace \texttt{PcfHigherCaso}), and \texttt{HermitePade.lean} (the irrationality
pipeline and the concrete certificates, namespace \texttt{PcfHermitePade}) are mirrored
in this catalogue under \texttt{verify/}.

\begin{remark}[Index convention]\label{rem:index}
The formalization parametrizes the order as $k+1$ via \texttt{Fin (k+1)} (so the
recurrence is $s(n+(k+1))=\sum_{j\le k}c_j(n)s(n+j)$ and the matrix is
$(k+1)\times(k+1)$); its sign therefore reads $(-1)^{k}$. This is the prose order-$k$
law~\eqref{eq:step} under $k\mapsto k+1$: an index shift only, with $(-1)^{(k+1)-1}
=(-1)^k$.
\end{remark}

\begin{theorem}[Formalized core; \textsc{proven}]\label{thm:lean}
The following twenty-six declarations are machine-checked with clean axiom cones ---
each a subset of $\{\texttt{propext},\texttt{Classical.choice},\texttt{Quot.sound}\}$, no
\texttt{sorryAx}, sources free of \texttt{sorry}/\texttt{admit}/\texttt{native\_decide},
\texttt{lake build} exit $0$.
\begin{itemize}\sloppy
\item \emph{Uniform-in-$k$ law} (\texttt{PcfGeneralCaso}, \texttt{GeneralCaso.lean}):
\texttt{casoMat\_det\_step} is the step law~\eqref{eq:step}
$\Caso(n+1)=(-1)^{k}c_0(n)\Caso(n)$ (order $k+1$; Remark~\ref{rem:index});
\texttt{casoMat\_det\_eq} is the closed product form~\eqref{eq:prod};
\texttt{casoMat\_det\_eq\_const} is the constant-coefficient power form
(Corollary~\ref{cor:const}); \texttt{casoMat\_det\_ne\_zero\_of\_init} is the
integral-domain non-vanishing certificate (Corollary~\ref{cor:indep}).
\item \emph{Fundamental-system converse} (\texttt{PcfGeneralCaso},
\texttt{GeneralCaso.lean}): \texttt{sol\_unique\_of\_init} (forward determinacy) and
\texttt{solution\_isLinearCombo} (completeness) --- Theorem~\ref{thm:converse}.
\item \emph{Low-order witnesses and the bridge} (\texttt{PcfHigherCaso},
\texttt{HigherCaso.lean}):
\texttt{caso2\_step} (the order-$2$ sign $-c_0$), \texttt{caso3\_step} (the order-$3$
step $+c_0$), \texttt{caso3\_eq} (its product form), \texttt{caso3\_eq\_const} (its
constant-coefficient case), \texttt{caso3\_ne\_zero\_of\_init} (its domain
non-vanishing certificate), and \texttt{caso3\_matches\_general} (the bespoke
$3\times3$ expansion equals $\det\CasoM$ of the general \texttt{casoMat} at $k=2$,
reconciling the witnesses with the general theorem) --- Proposition~\ref{prop:lowk}.
\item \emph{Irrationality pipeline} (\texttt{PcfHermitePade}, \texttt{HermitePade.lean}):
\texttt{caso\_mulVec\_apply} (window $=\CasoM(n)\,a$, \eqref{eq:window}),
\texttt{linForm\_window\_ne\_zero} (no full-window vanishing over a domain),
\texttt{linForm\_frequently\_ne\_zero} (frequent non-vanishing),
\texttt{irrational\_of\_linearForm\_frequently} (the linear-form irrationality
criterion), and the
capstone \texttt{irrational\_of\_fundamentalSet\_form} --- Section~\ref{sec:cert}.
\item \emph{Concrete order-$3$ certificates} (\texttt{PcfHermitePade},
\texttt{HermitePade.lean}): the geometric witness \texttt{gsol\_rec},
\texttt{gcaso0}, \texttt{gcaso\_ne\_zero} ($\Caso(n)=2\cdot6^n$),
\texttt{geom\_window\_ne\_zero}; and the golden-ratio certificate \texttt{fibSol\_rec},
\texttt{fibCaso0} (Cassini $\Caso(0)=1$), \texttt{fibCaso\_ne\_zero},
\texttt{fib\_remainder\_tendsto} ($F_n\varphi-F_{n+1}=-\psi^n\to0$), and
\texttt{goldenRatio\_irrational\_via\_caso} ($\varphi$ irrational).
\end{itemize}
Of the twenty-six, exactly two close with the minimal cone
$\{\texttt{propext},\texttt{Quot.sound}\}$ --- the by-hand witnesses
\texttt{caso2\_step} and \texttt{caso3\_step}; the remaining twenty-four additionally
use \texttt{Classical.choice}. This \texttt{Classical.choice} usage is inherited from
Mathlib's general infrastructure --- the determinant and permutation machinery
(\texttt{Matrix.det}) throughout (including the determinant side of the bridge
\texttt{caso3\_matches\_general}), the nonsingular matrix inverse
(\texttt{Matrix.mul\_nonsing\_inv}) in the converse, and the real-analysis limit layer
in the certificate --- not a sign that any proof is in a meaningful sense less
constructive than the by-hand cases; we report the cones factually rather than rank
them. No proof introduces any axiom outside the three-element clean set.
\end{theorem}

This follows the program's pattern: the law and converse have no labelled analytic
hypothesis at all, being finitary algebra over a \texttt{CommRing} (the domain and
completeness results add only \texttt{IsDomain\,R}, resp.\ a unit Casoratian); the
irrationality certificate isolates its \emph{single} analytic input --- the limit
$q(n)\alpha-p(n)\to0$ --- as a labelled hypothesis, discharged for the golden ratio by
Mathlib's Binet identity. The Casoratian
is defined as a genuine \texttt{Matrix} determinant, so the result is the honest
determinantal identity, not a bespoke $k=2,3$ polynomial expression --- and the order-$3$
by-hand expansion is itself certified equal to that determinant by
\texttt{caso3\_matches\_general} (Proposition~\ref{prop:lowk}). Only the
order-$2$ \emph{step} sign is recorded here, as \texttt{caso2\_step}: the full
order-$2$ closed-form suite already exists in \texttt{pcf-delta} as
\texttt{casoratian\_eqG}~\cite{pcfdelta} and is not duplicated, whereas order $3$
carries the full suite as the smallest genuinely higher-order witness.

\section{Independent symbolic verification}\label{sec:gate}

In keeping with the catalogue's gate-before-Lean discipline, a \texttt{sympy} harness
(\texttt{harness\_caso\_k.py}) re-derives the law independently of the Lean proof:
\begin{itemize}
\item a fully \emph{symbolic} proof at $k=3$: with generic symbolic data and generic
$c_0,c_1,c_2$, $\Caso(1)-c_0\,\Caso(0)$ expands to $0$ and $\Caso(1)$ is
free of $c_1,c_2$ (only $c_0$ enters);
\item \emph{exact-rational} checks of both the step law~\eqref{eq:step} and the
product form~\eqref{eq:prod} for $k=2,\dots,7$, over integer-coefficient recurrences
and initial data drawn from a \emph{fixed deterministic seed}
(\texttt{random.Random(7k+101t+5)} for trial $t$), so the runs are exactly
reproducible, confirming the alternating sign $(-1)^{k-1}$; and
\item named instances: the Tribonacci recurrence $s(n+3)=s(n+2)+s(n+1)+s(n)$ ($c_0\equiv1$)
has a \emph{constant} Casoratian, and, at $k=3$ (where the sign $(-1)^{k-1}=+1$), the
recurrence with $c_0(n)=n+1$ gives $\Caso(n)=n!\,\Caso(0)$ --- both
matching~\eqref{eq:prod} exactly.
\end{itemize}
All checks pass. The harness and the Lean sources are listed in
Section~\ref{sec:code}.

\section{Epistemic status}\label{sec:status}

\begin{itemize}
\item \textbf{PROVEN} (Lean, clean cones; Theorem~\ref{thm:lean}): the uniform-in-$k$
step law and product form, the constant-coefficient power form, the integral-domain
non-vanishing certificate, and the order-$2$/order-$3$ witnesses --- no analytic input.
\item \textbf{PROVEN} (Lean, clean cones; Theorem~\ref{thm:converse}): the
\emph{fundamental-system converse} --- forward determinacy
(\texttt{sol\_unique\_of\_init}) with no invertibility hypothesis, and completeness
(\texttt{solution\_isLinearCombo}) once $\Caso(0)$ is a unit. This was recorded as
\textsc{structural} in v1.0 and is now formalized in the uniform-in-$k$, time-varying
setting. The formalization in fact \emph{simplifies} the hypotheses we had anticipated:
because forward determinacy needs no reversibility, completeness follows from the
single assumption that $\Caso(0)$ is a unit, with no separate hypothesis that $c_0$ be a
unit. (Over a field this is the textbook fundamental-set theorem,
\cite[Ch.~2]{elaydi}; Mathlib's \texttt{LinearRecurrence} encodes only the
constant-coefficient case~\cite{mathlib}.)
\item \textbf{PROVEN} (Lean, clean cones; Section~\ref{sec:cert}): the
\emph{irrationality pipeline} --- the window identity~\eqref{eq:window}, the
domain-window non-vanishing, the linear-form irrationality criterion, and the capstone
\texttt{irrational\_of\_fundamentalSet\_form} --- together with its end-to-end
instantiation \texttt{goldenRatio\_irrational\_via\_caso} re-deriving $\varphi$
irrational through the order-$3$ Casoratian. The certificate carries exactly one
labelled analytic hypothesis (the remainder's limit), discharged here by Mathlib's
Binet identity. \emph{We claim no novelty for the irrationality of $\varphi$ itself}
(classical; already in Mathlib via $\sqrt5$): this is an independent Casoratian route
and a demonstration witness that the pipeline closes an irrationality theorem, in the
spirit of the by-hand faithfulness witnesses (Proposition~\ref{prop:lowk}), not a new
number-theoretic result --- see Remark~\ref{rem:honest-phi}.
\item \textbf{VERIFIED} (numerically/symbolically; Section~\ref{sec:gate}): the law at
$k=3$ symbolically and for $k=2,\dots,7$ by exact-rational checks, plus the named
instances.
\item \textbf{CONJECTURED / open}: a concrete simultaneous-Hermite--Pad\'e PCF whose
order-$3$ Casoratian drives the now-proven pipeline (Section~\ref{sec:cert}) to the
irrationality of a constant \emph{not already known} to be irrational --- an
Ap\'ery-style certificate ($\zeta(3)$-type) rather than the golden-ratio demonstration.
This is the genuinely open next entry: the determinantal machinery, the converse, and
the certificate engine it consumes are now all in place; what remains is to exhibit a
target recurrence whose shrinking remainder yields a \emph{new} irrationality theorem.
\end{itemize}
We claim no novelty for Abel's formula itself, which is classical
(\cite[Ch.~2]{elaydi},\cite{kelleypeterson},\cite{milnethomson}), nor for the
irrationality of the golden ratio; the contributions are the
uniform-in-$k$ machine-checked proof, the formalized fundamental-system converse, the
reusable machine-checked irrationality pipeline, and the opening of the order-$\ge3$
direction in this corpus.

\section{Relation to prior work and novelty}\label{sec:related}

\paragraph{Formal-methods landscape.}
The two standard formal-library treatments of linear recurrences both target the
\emph{constant-coefficient} solution theory, not the Casoratian evolution law.
Mathlib's \texttt{LinearRecurrence}~\cite{mathlib} fixes an order and a tuple of
\emph{constant} coefficients \texttt{coeffs : Fin order $\to$ R}, builds the solution
submodule of $\N\to R$, the \texttt{LinearEquiv} sending a solution to its first
\texttt{order} terms (so solutions are pinned by their initial data), and the
characteristic polynomial; it defines no Casoratian and proves no Abel-type evolution
law. Eberl's \emph{Linear Recurrences} entry in the Isabelle Archive of Formal
Proofs~\cite{eberl} develops constant-coefficient recurrences via formal power series
and polynomial factorisation into an executable closed-form solver, again without the
Casoratian/Abel identity. Our object is complementary and, in one direction, strictly
more general: the coefficients $c_j(n)$ are arbitrary functions of $n$ (time-varying),
the Casoratian is a genuine \texttt{Matrix} determinant, and the theorem is the
determinantal evolution law uniform in the order $k$ --- exactly the
non-vanishing / linear-independence certificate (Corollary~\ref{cor:indep}) that
Hermite--Pad\'e and Ap\'ery-type constructions consume (\cite{vdpoorten,beukers}). A
machine-checked Casoratian does already exist for the order-$2$ case in another system:
the Coq/MathComp formal proof of the irrationality of $\zeta(3)$~\cite{apery} defines
the $2\times2$ Wronskian \texttt{ba\_casoratian} of the two Ap\'ery sequences and
derives its first-order (Abel) law, but only for that concrete instance --- indeed its
source remarks that the order-$1$ recurrence ``should be obtained from something more
general.'' Having searched Mathlib, the Isabelle AFP, and the Coq/MathComp ecosystem by
name (\emph{Casoratian}, \emph{Wronskian}) and by structure (as of June~2026), we find no
\emph{machine-checked} general-order, uniform-in-$k$ Casoratian Abel formula with
time-varying coefficients in any of these systems --- precisely the ``something more
general'' the Ap\'ery formalization lacked. The \emph{identity} itself is of course
classical and not in question: the general-order law with time-varying coefficients is
textbook (\cite[Ch.~2]{elaydi},\cite{kelleypeterson},\cite{milnethomson}), for which we
claim no novelty. Our claim is strictly one of \emph{formalization coverage} --- the
machine-checked theorem, uniform in $k$ and with time-varying $c_j(n)$, that the present
note contributes and gate-checks. The irrationality pipeline of
Section~\ref{sec:cert} is likewise complementary to~\cite{apery}: where that development
hard-wires a $2\times2$ Casoratian into a single bespoke $\zeta(3)$ certificate, ours is
a \emph{reusable} order-$(k{+}1)$ engine that consumes any integer fundamental set and a
shrinking real remainder, here demonstrated on the classical golden ratio (for which we
claim no number-theoretic novelty, Remark~\ref{rem:honest-phi}) and awaiting a target
not already known irrational (Section~\ref{sec:status}).

\paragraph{Within this corpus, and a note on hosting.}
The deposited corpus had formalized the Casoratian only in its order-$2$ form --- the
$4/\pi$ PCF Casoratian closed form (Theorem~6.3 of the ladder paper~\cite{ladder};
machine-verified in~\cite{caso2lean}) and the general order-$2$
\texttt{casoratian\_eqG} in \texttt{pcf-delta}~\cite{pcfdelta} --- so this is the
program's first order-$\ge3$ result and the backbone of a planned Hermite--Pad\'e /
Ap\'ery line. We deposit it as a standalone catalogue entry rather than as a Mathlib
pull request for two reasons: the result is more general than Mathlib's
constant-coefficient \texttt{LinearRecurrence} API and would need new
time-varying-coefficient scaffolding to land there, and the corpus requires the core
to be hosted and cone-gated inside \texttt{pcf-delta}'s \texttt{Check.lean}. A Mathlib
contribution --- a \texttt{Casoratian} determinant and its Abel evolution law,
specialising to \texttt{LinearRecurrence} in the constant-coefficient case --- is a
natural next step and would convert ``not otherwise available'' into an externally
reviewed library lemma. For the record, the Lean and \texttt{sympy} sources label this
identity ``Abel--Jacobi--Liouville''; the standard discrete name, used here, is Abel's
formula (Elaydi~\cite{elaydi}), the triple form naming the continuous-Wronskian
analogue (Liouville / Jacobi). The note supplements the SIARC program~\cite{siarc}.

\section{Code and data availability}\label{sec:code}

\begin{sloppypar}
The Lean sources are \texttt{verify/GeneralCaso.lean} (namespace
\texttt{PcfGeneralCaso}; the uniform-in-$k$ law and the fundamental-system converse),
\texttt{verify/HigherCaso.lean} (namespace \texttt{PcfHigherCaso}; the by-hand
low-order witnesses), and \texttt{verify/HermitePade.lean} (namespace
\texttt{PcfHermitePade}; the irrationality pipeline and the concrete order-$3$
certificates); they are built and axiom-cone-checked inside the
\texttt{pcf-delta} \texttt{PcfContinuant} project~\cite{pcfdelta} via
\texttt{lake env lean Check.lean} (toolchain \texttt{leanprover/lean4:v4.30.0},
Mathlib \texttt{rev=v4.30.0}); the verbatim \texttt{\#print axioms} output for the
twenty-six declarations is mirrored in this repository as
\texttt{verify/AXIOM\_CONES.txt}, so the cone partition of Section~\ref{sec:lean} is
reproducible here without rebuilding the host project. The independent symbolic gate is
\texttt{harness\_caso\_k.py} (\texttt{sympy}~\cite{sympy}; run
\texttt{python harness\_caso\_k.py}), which prints the $k=3$ symbolic proof, the
exact-rational $k=2,\dots,7$ verification, and the named instances. Public repository:
\url{https://github.com/papanokechi/pcf-casoratian-catalogue}.
\end{sloppypar}

\begin{thebibliography}{99}

\bibitem{elaydi}
S.~Elaydi,
\emph{An Introduction to Difference Equations}, 3rd ed.,
Undergraduate Texts in Mathematics, Springer, 2005.
The Casoratian and Abel's lemma for its evolution under variable coefficients are
treated in Ch.~2 (linear difference equations; linear independence of solutions).

\bibitem{kelleypeterson}
W.~G.~Kelley and A.~C.~Peterson,
\emph{Difference Equations: An Introduction with Applications},
Academic Press, 1991. Casoratian (discrete Wronskian) and the linear-independence
theory for linear difference equations.

\bibitem{milnethomson}
L.~M.~Milne-Thomson,
\emph{The Calculus of Finite Differences},
Macmillan, 1933 (reprinted, AMS Chelsea Publishing).
A classical treatment of the Casoratian and linear difference equations.

\bibitem{vdpoorten}
A.~J.~van der Poorten,
\emph{A proof that Euler missed\ldots\ Ap\'ery's proof of the irrationality of
$\zeta(3)$},
The Mathematical Intelligencer \textbf{1}(4) (1979), 195--203.
DOI \texttt{10.1007/BF03028234} (\url{https://doi.org/10.1007/BF03028234}).
The Hermite--Pad\'e / simultaneous-approximation context in which higher-order
Casoratians arise.

\bibitem{beukers}
F.~Beukers,
\emph{A note on the irrationality of $\zeta(2)$ and $\zeta(3)$},
Bull.\ London Math.\ Soc.\ \textbf{11}(3) (1979), 268--272.
DOI \texttt{10.1112/blms/11.3.268} (\url{https://doi.org/10.1112/blms/11.3.268}).
A Pad\'e / Legendre-polynomial route to Ap\'ery's irrationality theorems --- the
higher-order-recurrence setting in which Casoratians certify independence.

\bibitem{mathlib}
The mathlib Community,
\emph{The Lean Mathematical Library},
in: Proc.\ 9th ACM SIGPLAN Int.\ Conf.\ on Certified Programs and Proofs (CPP~2020),
ACM, 2020, 367--381. DOI \texttt{10.1145/3372885.3373824}
(\url{https://doi.org/10.1145/3372885.3373824}).

\bibitem{lean4}
L.~de~Moura and S.~Ullrich,
\emph{The Lean~4 Theorem Prover and Programming Language},
in: Automated Deduction --- CADE~28, Lecture Notes in Computer Science \textbf{12699},
Springer, 2021, 625--635. DOI \texttt{10.1007/978-3-030-79876-5\_37}
(\url{https://doi.org/10.1007/978-3-030-79876-5_37}).

\bibitem{sympy}
A.~Meurer et al.,
\emph{SymPy: symbolic computing in Python},
PeerJ Computer Science \textbf{3} (2017), e103.
DOI \texttt{10.7717/peerj-cs.103} (\url{https://doi.org/10.7717/peerj-cs.103}).

\bibitem{eberl}
M.~Eberl,
\emph{Linear Recurrences},
Archive of Formal Proofs, October 2017.
\url{https://www.isa-afp.org/entries/Linear_Recurrences.html}.
An Isabelle/HOL formalization of constant-coefficient linear recurrences and an
executable closed-form solver (no Casoratian / Abel evolution law).

\bibitem{apery}
F.~Chyzak, A.~Mahboubi, T.~Sibut-Pinote, and E.~Tassi,
\emph{A Computer-Algebra-Based Formal Proof of the Irrationality of $\zeta(3)$},
in: Interactive Theorem Proving (ITP~2014), Lecture Notes in Computer Science,
Springer, 2014, 160--176. DOI \texttt{10.1007/978-3-319-08970-6\_11}
(\url{https://doi.org/10.1007/978-3-319-08970-6_11}).
The Coq/MathComp development (\texttt{coq-community/apery}) defines a concrete
order-$2$ Casoratian \texttt{ba\_casoratian} and its first-order law.

\bibitem{caso2lean}
Papanokechi,
\emph{Machine-verified Lean~4/Mathlib formalization of the Casoratian closed form
(Theorem~6.3) for a $4/\pi$ polynomial continued fraction}, 2026.
The order-$2$ predecessor this note extends. Zenodo concept DOI
\texttt{10.5281/zenodo.20500490} (\url{https://doi.org/10.5281/zenodo.20500490}).

\bibitem{ladder}
Papanokechi,
\emph{Polynomial Continued Fractions: a Proved Logarithmic Ladder, a $4/\pi$
Casoratian Identity, and 482 Irrational Constants}, 2026.
Zenodo concept DOI \texttt{10.5281/zenodo.19491767}
(\url{https://doi.org/10.5281/zenodo.19491767}).

\bibitem{pcfdelta}
Papanokechi,
\emph{pcf-delta: the growth-law correction constant $\delta=\log R_\infty$ of the
quadratic polynomial continued fraction}, 2026.
Hosts and builds the Lean \texttt{PcfContinuant} project and \texttt{Check.lean} gate
in which the present core is cone-checked. Zenodo concept DOI
\texttt{10.5281/zenodo.20578400} (\url{https://doi.org/10.5281/zenodo.20578400}).

\bibitem{siarc}
Papanokechi,
\emph{SIARC: a program statement for structured, incrementally-audited research on
continued fractions}, 2026.
Zenodo concept DOI \texttt{10.5281/zenodo.19885549}
(\url{https://doi.org/10.5281/zenodo.19885549}).

\end{thebibliography}

\paragraph{AI-assistance disclosure.}
Prepared with AI assistance (GitHub Copilot; Anthropic Claude) for drafting,
computer-algebra verification, and manuscript preparation. All symbolic and
exact-rational checks were obtained by running the scripts named in the Code and Data
Availability section; the Lean axiom cones were obtained by \texttt{\#print axioms}. The
author takes responsibility for the content and for the placement of each claim within
the four-grade epistemic partition. No AI system is credited as an author, and all
AI-assisted text and computations were checked by the author.

\end{document}
